%%% -*- coding: utf-8 -*-
% !TeX root = main-slides.tex

\lecture{C\#}{csharp-basics}
\section{C\#}

\showtitlepage{C\# Basics for Java Programmers}

\subsection{Java vs.\ C\#}

\begin{frame}{Java vs.\ C\#: Big Picture}
  \begin{itemize}
    \item Both are class-based, object-oriented, statically typed languages
      with a similar-looking syntax (curly braces, semicolons, single
      inheritance of classes, interfaces).
    \item Both compile to an intermediate bytecode that runs on a virtual
      machine with a JIT compiler and garbage collection.
      \begin{itemize}
        \item Java: source $\to$ \texttt{.class} bytecode $\to$ JVM
        \item C\#: source $\to$ intermediate language (\texttt{.dll}/\texttt{.exe})
          $\to$ CLR (in Unity: Mono or IL2CPP)
      \end{itemize}
    \item Despite the similarity, several design decisions differ --
      this section highlights the ones you will notice first.
  \end{itemize}
\end{frame}

\begin{frame}{Naming Conventions}
  \begin{itemize}
    \item Java: \texttt{camelCase} for methods and fields,
      \texttt{PascalCase} for classes.
    \item C\#: \texttt{PascalCase} for classes, methods, \emph{and}
      properties; \texttt{camelCase} for local variables and parameters.
    \item Consequence: a method that reads a value is
      \texttt{getScore()} in Java, but simply \texttt{Score} (a property,
      see below) or \texttt{GetScore()} in C\#.
  \end{itemize}
\end{frame}

\begin{frame}[fragile]{Properties Instead of Getters/Setters}
  \begin{columns}[T]
    \column{0.5\textwidth}
      \textbf{Java}
      \begin{lstlisting}[language=Java]
private int score;

public int getScore() {
  return score;
}
public void setScore(int s) {
  score = s;
}
      \end{lstlisting}
    \column{0.5\textwidth}
      \textbf{C\#}
      \begin{lstlisting}
private int score;

public int Score {
  get { return score; }
  set { score = value; }
}

// or, auto-implemented:
public int Score { get; set; }
      \end{lstlisting}
  \end{columns}
  Properties look like fields to the caller (\texttt{obj.Score = 3;})
  but behave like methods.
\end{frame}

\begin{exerciseframe}{Port a Java Bean}
  Rewrite this class in C\#, following the C\# naming conventions and
  replacing the getter/setter pair by a property. Adapt the caller as well.
  \vspace{1ex}

  \begin{columns}[T]
    \column{0.52\textwidth}
      \textbf{Counter.java}
      \begin{lstlisting}[language=Java,basicstyle=\scriptsize\ttfamily]
public class Counter {
  private int count;

  public int getCount() {
    return count;
  }
  public void setCount(int c) {
    count = c;
  }
  public void increment() {
    count++;
  }
}
      \end{lstlisting}
    \column{0.48\textwidth}
      \textbf{Caller}
      \begin{lstlisting}[language=Java,basicstyle=\scriptsize\ttfamily]
Counter c = new Counter();
c.setCount(5);
c.increment();
System.out.println(c.getCount());
      \end{lstlisting}
  \end{columns}
\end{exerciseframe}

\begin{solutionframe}{Port a Java Bean}
  \begin{columns}[T]
    \column{0.52\textwidth}
      \textbf{Counter.cs}
      \begin{lstlisting}[basicstyle=\scriptsize\ttfamily]
public class Counter
{
  public int Count { get; set; }

  public void Increment()
  {
    Count++;
  }
}
      \end{lstlisting}
    \column{0.48\textwidth}
      \textbf{Caller}
      \begin{lstlisting}[basicstyle=\scriptsize\ttfamily]
Counter c = new Counter();
c.Count = 5;
c.Increment();
Console.WriteLine(c.Count);
      \end{lstlisting}
  \end{columns}

  \vspace{3ex}

  \hint{The auto-implemented property provides the backing field. Properties
  and methods are \texttt{PascalCase}, while the local variable \texttt{c}
  stays \texttt{camelCase}.}
\end{solutionframe}

\begin{exerciseframe}{Validation Without Breaking Callers}
  \texttt{Count} shall reject negative values by throwing an
  \texttt{ArgumentOutOfRangeException}. Change \texttt{Counter} only -- this
  call site must keep compiling unchanged.
  \begin{lstlisting}
Counter c = new Counter();
c.Count = 5;    // still fine
c.Count = -1;   // must throw
  \end{lstlisting}
  Which change would a Java program with \texttt{getCount}/\texttt{setCount}
  have required here?
\end{exerciseframe}

\begin{solutionframe}{Validation Without Breaking Callers}
  \begin{lstlisting}[basicstyle=\scriptsize\ttfamily]
private int count;
public int Count
{
  get => count;
  set
  {
    if (value < 0)
    {
      throw new ArgumentOutOfRangeException(nameof(value));
    }
    count = value;
  }
}
  \end{lstlisting}

  \hint{\texttt{value} is the implicit parameter of the setter, and
  \texttt{nameof(value)} turns its name into the string \texttt{"value"}. Since
  call sites never change, you can start with an auto-implemented property and
  add validation later -- which is why C\# exposes properties where Java
  exposes getters and setters.}
\end{solutionframe}

\begin{exerciseframe}{A Computed Property}
  Add a read-only \texttt{IsEmpty} to \texttt{Counter} that yields
  \texttt{true} exactly when \texttt{Count} is zero, without storing any
  additional state.
  \begin{lstlisting}
Counter c = new Counter();
Console.WriteLine(c.IsEmpty);   // True
c.Count = 3;
Console.WriteLine(c.IsEmpty);   // False
  \end{lstlisting}
  When would you prefer a method \texttt{bool IsEmpty()} over a property?
\end{exerciseframe}

\begin{solutionframe}{A Computed Property}
  \begin{columns}[T]
    \column{0.5\textwidth}
      \textbf{Expression-bodied}
      \begin{lstlisting}
public bool IsEmpty
  => Count == 0;
      \end{lstlisting}
    \column{0.5\textwidth}
      \textbf{Equivalent long form}
      \begin{lstlisting}
public bool IsEmpty
{
  get { return Count == 0; }
}
      \end{lstlisting}
  \end{columns}
  A property fits when the value is cheap to compute, free of side effects and
  reads like state. Prefer a method when the call does real work, may fail or
  needs parameters.
  \vspace{3ex}

  \hint{Omitting \texttt{set} makes the property read-only. Callers cannot
  tell \texttt{IsEmpty} from a field.}
\end{solutionframe}

\begin{frame}{Value Types vs.\ Reference Types}
  \framelabel{slide:valuetypes}
  \begin{itemize}
    \item Java: primitives (\texttt{int}, \texttt{double}, \texttt{boolean},
      \ldots) are value types; \emph{everything else} is a reference type
      and lives on the heap.
    \item C\#: you can define your own value types with \texttt{struct}
      (copied by value, usually on the stack) in addition to
      \texttt{class} (reference type, on the heap).
    \item Relevant for Unity: \texttt{Vector3}, \texttt{Quaternion},
      \texttt{Color} are \texttt{struct}s -- assigning or passing them
      copies the whole value. A common bug when coming from Java is
      expecting mutations to a struct field of another object to
      "stick".
    \item Whenever a property, an indexer or a method returns a struct, it
      returns a copy. Assigning to a field of that copy is rejected by the
      compiler (error CS1612), as the copy would be discarded right away.
  \end{itemize}
\end{frame}

\begin{exerciseframe}{Predict the Output}
  What does this program print? Write down your answer \emph{before} you run it.
  \vspace{1ex}

  \begin{lstlisting}[basicstyle=\scriptsize\ttfamily]
struct Point { public int X; }

class Holder
{
  public Point P = new Point();
}

Holder h = new Holder();
Point copy = h.P;
copy.X = 42;
Console.WriteLine(h.P.X);
  \end{lstlisting}
  \vspace{1ex}

  Then replace \texttt{struct Point} by \texttt{class Point} and predict again.
\end{exerciseframe}

\begin{solutionframe}{Predict the Output}
  \begin{itemize}
    \item \texttt{struct Point} prints \texttt{0}:\\
       \texttt{Point copy = h.P;} copies the value,\\
       hence \texttt{copy.X = 42;} leaves \texttt{h.P}
       untouched.
    \item \texttt{class Point} prints \texttt{42}:\\
     \texttt{copy} and \texttt{h.P} now refer to the same object.
  \end{itemize}
  \vspace{3ex}

  \hint{Java offers only the second case, as every user-defined type is a
  reference type there. The first case is what makes Unity's \texttt{Vector3},
  \texttt{Quaternion} and \texttt{Color} behave the way they do.}
\end{solutionframe}

\begin{exerciseframe}{Why the Compiler Refuses}
  With \texttt{Point} declared as a \texttt{struct}, this does not compile.\\
  Read the error message, then make it work in two different ways.
  \vspace{1ex}

  \begin{lstlisting}[basicstyle=\scriptsize\ttfamily]
List<Point> points = new List<Point> { new Point() };
points[0].X = 5;
  \end{lstlisting}
  \vspace{1ex}

  Afterwards replace \texttt{List<Point>} by \texttt{Point[]} and explain why
  \emph{that} version compiles.
\end{exerciseframe}

\begin{solutionframe}{Why the Compiler Refuses}
  The compiler reports the C\# error CS1612.
  \vspace{1ex}

  The indexer of \texttt{List<T>} is a method returning a \emph{copy} of the
  element,\\
  so \texttt{points[0].X = 5;} would modify a temporary that is
  discarded immediately.\\
  \vspace{1ex}
  \pause

  \textbf{First fix:}\\
  \begin{lstlisting}[basicstyle=\scriptsize\ttfamily]
Point p = points[0];   // read
p.X = 5;               // modify
points[0] = p;         // write back
  \end{lstlisting}
  \pause

  \textbf{Second fix:}\\
  Declaring \texttt{Point} as a class is the other fix: the indexer then hands
  out a reference and the original assignment compiles.
  \vspace{3ex}
  \pause

  \hint{An array behaves differently: \texttt{array[0]} denotes a variable
  rather than a method call, hence \texttt{array[0].X = 5;} compiles even for a
  struct.}
\end{solutionframe}

\begin{exerciseframe}{Mutating a Unity Vector3}
  A classic Unity compile error. Explain it, then fix it in two ways.
  \vspace{1ex}

  \begin{lstlisting}[basicstyle=\scriptsize\ttfamily]
transform.position.x = 5f;
  \end{lstlisting}
  And why does \texttt{transform.position += Vector3.up;} compile?
\end{exerciseframe}

\begin{solutionframe}{Mutating a Unity Vector3}
  \texttt{position} is a property returning a \texttt{Vector3}, and
  \texttt{Vector3} is a \texttt{struct}. Its getter hands out a copy, so the
  assignment would modify a temporary -- the same error CS1612 as on
  \slideref{slide:valuetypes}.
  \vspace{1ex}

\textbf{First fix:}\\
  \begin{lstlisting}[basicstyle=\scriptsize\ttfamily]
Vector3 p = transform.position;
p.x = 5f;
transform.position = p;
  \end{lstlisting}

  \textbf{Second fix:}\\
 \begin{lstlisting}[basicstyle=\scriptsize\ttfamily]
transform.position =
  new Vector3(5f, transform.position.y, transform.position.z);
  \end{lstlisting}
  \vspace{3ex}

  \hint{\texttt{transform.position += Vector3.up;} works because it expands
  into a read, an addition and an assignment back through the setter --
  precisely the pattern above.}
\end{solutionframe}

\begin{frame}{Nullability}
  \begin{itemize}
    \item Java: any reference may be \texttt{null}; primitives cannot be
      \texttt{null}.
    \item C\#: reference types may be \texttt{null} just like in Java;
      value types are \emph{not} nullable unless explicitly wrapped,
      e.g.\ \texttt{int?\ x = null;} (short for \texttt{Nullable<int>}).
    \item \texttt{Nullable<T>} provides \texttt{HasValue} to ask whether a value
      is present and \texttt{Value} to unwrap it. Reading \texttt{Value} of an
      empty nullable throws an \texttt{InvalidOperationException}.
    \item Convenience operators without a Java equivalent:
      \begin{itemize}
        \item \texttt{a?.Foo()} -- null-conditional call, evaluates to
          \texttt{null} instead of throwing if \texttt{a} is \texttt{null}.
        \item \texttt{a ?? b} -- null-coalescing, yields \texttt{b} if
          \texttt{a} is \texttt{null}.
      \end{itemize}
  \end{itemize}
\end{frame}

\begin{exerciseframe}{A Nullable Return Value}
  Implement \texttt{SafeDivide} so that it yields \texttt{null} instead of
  throwing when the divisor is zero. What do the three statements print?
  \vspace{1ex}

  \begin{lstlisting}[basicstyle=\scriptsize\ttfamily]
static int? SafeDivide(int a, int b)
{
  // your code here
}

Console.WriteLine(SafeDivide(10, 2) ?? -1);
Console.WriteLine(SafeDivide(10, 0) ?? -1);
Console.WriteLine(SafeDivide(10, 0).Value);
  \end{lstlisting}
  \vspace{1ex}

  Why can the return type not simply be \texttt{int}?
\end{exerciseframe}

\begin{solutionframe}{A Nullable Return Value}
  \begin{lstlisting}[basicstyle=\scriptsize\ttfamily]
static int? SafeDivide(int a, int b)
{
  if (b == 0)
  {
    return null;
  }
  return a / b;
}
  \end{lstlisting}
  The first and second statements print \texttt{5} and \texttt{-1}, respectively, and the third one
  throws an \texttt{InvalidOperationException}. An \texttt{int} cannot be
  \texttt{null}, hence the result has to be wrapped.
  \vspace{1ex}
  \vspace{3ex}

  \hint{\texttt{int?} is \texttt{Nullable<int>}, a struct offering
  \texttt{HasValue} and \texttt{Value}. Query \texttt{HasValue} or supply a
  fallback with \texttt{??} rather than reaching for \texttt{Value}. Java would
  need \texttt{Integer} or \texttt{Optional<Integer>} here.}
\end{solutionframe}

\begin{exerciseframe}{Chaining the Null-Conditional Operator}
  Compute the length of the ZIP code in a single expression, yielding
  \texttt{-1} if the person, the address or the ZIP code is \texttt{null}.
  \vspace{1ex}

  \begin{lstlisting}[basicstyle=\scriptsize\ttfamily]
class Address { public string Zip { get; set; } }
class Person { public Address Home { get; set; } }

Person person = null;
int length = // your expression here
  \end{lstlisting}
  \vspace{1ex}

  What is the static type of \texttt{person?.Home?.Zip?.Length}?
\end{exerciseframe}

\begin{solutionframe}{Chaining the Null-Conditional Operator}
  \begin{lstlisting}[basicstyle=\scriptsize\ttfamily]
int length = person?.Home?.Zip?.Length ?? -1;
  \end{lstlisting}
  \begin{itemize}
    \item The type of \texttt{person?.Home?.Zip?.Length} is \texttt{int?}, not
      \texttt{int}:\\
       \texttt{?.} lifts the result of a value type to its
      nullable counterpart.
    \item The chain short-circuits. As soon as one link is \texttt{null}, the
      whole expression is \texttt{null} and the remaining links are never
      evaluated.
  \end{itemize}
  \vspace{3ex}

  \hint{Java requires nested \texttt{if} statements or a chain of
  \texttt{Optional.map} calls for this.}
  \vspace{1ex}

  \hint{\texttt{?.} guards method calls as well, e.g.\ \texttt{list?.Clear()}.}
\end{solutionframe}

\begin{frame}{Namespaces and Packages}
  \begin{itemize}
    \item Java: package name must match the directory structure; one
      public top-level class per file; \texttt{import} for other packages.
    \item C\#: \texttt{namespace} is independent of the folder layout;
      a file may contain several classes and even several namespaces;
      \texttt{using} instead of \texttt{import}.
  \end{itemize}
\end{frame}

\begin{frame}{Assemblies}
  \begin{itemize}
    \item An assembly is the unit C\# code is compiled into: the \texttt{.dll}
      or \texttt{.exe} holding intermediate language and metadata. One project
      yields exactly one assembly.
    \item Assemblies and namespaces are independent of each other: an assembly
      may contain many namespaces, and a namespace may span several assemblies.
    \item Using a type from another assembly requires a \emph{reference} to
      that assembly. A \texttt{using} directive only abbreviates names, it does
      not make types available.
    \item The assembly is the boundary \texttt{internal} refers to. The Java
      counterparts are the JAR, which only packages, and the module of Java 9,
      which does restrict access across its boundary.
  \end{itemize}
  \vspace{1ex}

  \hint{In Unity each \texttt{.asmdef} file defines one assembly; scripts not
  covered by any of them land in \texttt{Assembly-CSharp}. As compilation
  happens per assembly, splitting a project into several of them also shortens
  recompile times.}
\end{frame}

\begin{frame}[fragile]{Partial Classes: One Type, Multiple Files}
  \framelabel{slide:partial}
  \begin{itemize}
    \item Java: a (non-nested) class must be fully defined in a single file.
    \item C\#: the \texttt{partial} keyword lets you split the definition
      of one class, struct, or interface across several files -- the
      compiler merges all parts into a single type.
  \end{itemize}
  \begin{columns}[T]
    \column{0.5\textwidth}
      \textbf{Player.cs}
      \begin{lstlisting}
public partial class Player
{
  public int Health;

  public void Move() { ... }
}
      \end{lstlisting}
    \column{0.5\textwidth}
      \textbf{Player.Combat.cs}
      \begin{lstlisting}
public partial class Player
{
  public void Attack() { ... }
}
      \end{lstlisting}
  \end{columns}

  Every part needs the \texttt{partial} keyword, lives in the
  same namespace and assembly, and must agree on access modifier
  and base class/interfaces.

  \hint{Can be used to split large classes into more cohesive parts.
  However, it is generally better to avoid larges classes altogether.}

\end{frame}

\begin{frame}{Type Names and File Names}
  \begin{itemize}
    \item Java: a public top-level class has to reside in a file of the very
      same name, and the compiler enforces that.
    \item C\#: no such rule. A file may contain any number of types, public or
      not, and its name need not relate to any of them. Conversely,
      \texttt{partial} lets a single type span several files.
    \item The convention is nevertheless one top-level type per file, named
      after that type, and \texttt{Type.Aspect.cs} for the parts of a partial
      type, as with \texttt{Player.Combat.cs} on \slideref{slide:partial}. Type
      parameters are
      dropped: \texttt{Graph<T>} lives in \texttt{Graph.cs}.
  \end{itemize}

  \hint{Unity turns that convention into a rule for \texttt{MonoBehaviour}
      and \texttt{ScriptableObject}: unless the file is named after the class,
      the script cannot be attached as a component.}
  \vspace{1ex}

  \hint{Let type and file name match even where the compiler does not care.
  Readers, IDE navigation, reviews and diffs all rely on finding \texttt{Foo}
  in \texttt{Foo.cs}.}
\end{frame}

\begin{frame}{Access Modifiers}
  \framelabel{slide:access}
  \begin{itemize}
    \item Java: \texttt{public}, \texttt{protected}, \texttt{private},
      and \emph{package-private} (the default -- no keyword).
    \item C\#: \texttt{public}, \texttt{protected}, \texttt{private},
      \texttt{internal} (visible within the whole assembly, not just a
      folder/package), plus combinations
      \texttt{protected internal} and \texttt{private protected} (see below).
    \item C\# has no direct equivalent of Java's package-private; the
      default access (with no modifier) is \texttt{private} for members
      and \texttt{internal} for top-level types.
  \end{itemize}
\end{frame}

\begin{frame}{Access Modifiers: The Two Combinations}
  \begin{itemize}
    \item \texttt{protected internal} means protected \emph{or} internal and is
      therefore \emph{wider} than either part alone: the whole assembly plus
      derived classes everywhere.
    \item \texttt{private protected} means protected \emph{and} internal and is
      therefore \emph{narrower} than either part alone: only derived classes
      within the same assembly.
  \end{itemize}
  \vspace{1ex}

  \begin{center}
    \footnotesize
    \begin{tabular}{lcccc}
      \hline
      \textbf{Modifier}
        & \shortstack{unrelated class,\\same assembly}
        & \shortstack{derived class,\\same assembly}
        & \shortstack{derived class,\\other assembly}
        & \shortstack{anywhere\\else}\\
      \hline
      \texttt{internal}           & yes & yes & no  & no\\
      \texttt{protected}          & no  & yes & yes & no\\
      \texttt{protected internal} & yes & yes & yes & no\\
      \texttt{private protected}  & no  & yes & no  & no\\
      \hline
    \end{tabular}
  \end{center}
  \vspace{1ex}

  \hint{The unit of accessibility is the assembly, in Unity determined by the
  \texttt{.asmdef} files. Without any, all scripts end up in
  \texttt{Assembly-CSharp}, where \texttt{internal} is effectively
  \texttt{public} and \texttt{private protected} collapses to
  \texttt{protected}.}
\end{frame}

\begin{exerciseframe}{One File, Two Namespaces}
  Put both classes into a \emph{single} file whose name matches neither of them,
  in two different namespaces, and let \texttt{NodeView} use \texttt{Node}.
  \vspace{1ex}

  \begin{lstlisting}[basicstyle=\scriptsize\ttfamily]
// file: totally-unrelated-name.cs
namespace Company.Model
{
  public class Node { public string Name { get; set; } }
}

namespace Company.View
{
  public class NodeView { /* use Node here */ }
}
  \end{lstlisting}
  \vspace{1ex}

  How many files, and which file names, would Java have required?
\end{exerciseframe}

\begin{solutionframe}{One File, Two Namespaces}

  \texttt{Company.View.NodeView} using  \texttt{Company.Model.Node}:

  \begin{lstlisting}[basicstyle=\scriptsize\ttfamily]
namespace Company.View
{
  using Company.Model;

  public class NodeView
  {
    private Node node = new Node();
  }
}
  \end{lstlisting}
  Java would have required two files, \texttt{Node.java} and
  \texttt{NodeView.java}, in directories mirroring the package names.
  \vspace{2ex}

  \hint{A \texttt{using} directive may sit at the top of the file or inside a
  namespace, where it then applies to that namespace only. Without it, the
  fully qualified \texttt{Company.Model.Node} does the job as well.}
\end{solutionframe}

\begin{exerciseframe}{Breaking the Partial Rules}
  Split \texttt{Player} across \texttt{Player.cs} and
  \texttt{Player.Combat.cs} as shown on \slideref{slide:partial}. Then provoke
  two compilation errors, one at a
  time, and predict the message before compiling.
  \vspace{1ex}

  \begin{itemize}
    \item Drop the \texttt{partial} keyword from one of the two declarations.
    \item Let the two parts declare \emph{different} base classes.
  \end{itemize}
  \vspace{1ex}

  And does every part have to name the base class?
\end{exerciseframe}

\begin{solutionframe}{Breaking the Partial Rules}
  \begin{itemize}
    \item Missing keyword -- \texttt{CS0260: Missing partial modifier on
      declaration of type 'Player'; another partial declaration of this type
      exists.}
    \item Different base classes -- \texttt{CS0263: Partial declarations of
      'Player' must not specify different base classes.}
    \item No. It suffices that \emph{one} part names the base class, the others
      may omit it. They must only never contradict each other.
  \end{itemize}
  \vspace{4ex}

  \hint{All parts are merged into a single type, hence any disagreement between
  them is an error.}
\end{solutionframe}

\begin{exerciseframe}{The Default Access Modifier}
  This does not compile. Explain why, using the rule from
  \slideref{slide:access},
  then give two different fixes.
  \vspace{1ex}

  \begin{lstlisting}[basicstyle=\scriptsize\ttfamily]
class Config
{
  public int Volume { get; set; }
}

public class Service
{
  public Config GetConfig()
  {
    return new Config();
  }
}
  \end{lstlisting}
\end{exerciseframe}

\begin{solutionframe}{The Default Access Modifier}
  \texttt{Config} carries no access modifier and is therefore
  \texttt{internal}, whereas \texttt{GetConfig} is \texttt{public}. A public
  member must not expose a type less accessible than itself:
  \vspace{1ex}

  \begin{lstlisting}[language={},basicstyle=\scriptsize\ttfamily]
CS0050: Inconsistent accessibility: return type
'Config' is less accessible than method
'Service.GetConfig()'
  \end{lstlisting}
  Either declare \texttt{Config} as \texttt{public}, or restrict
  \texttt{GetConfig} to \texttt{internal}.
  \vspace{4ex}

  \hint{Nested types default to \texttt{private} rather than
  \texttt{internal} -- the default depends on where the type is declared.}
\end{solutionframe}

\begin{frame}{Delegates, Lambdas, and Events}
  \begin{itemize}
    \item Java: callbacks are modeled with functional interfaces
      (\texttt{Runnable}, \texttt{Comparator<T>}, \texttt{@FunctionalInterface})
      plus lambda expressions.
    \item C\#: \texttt{delegate} is a built-in, type-safe function-pointer
      type; \texttt{Action}/\texttt{Func<...>} are predefined generic
      delegates so you rarely declare your own.
    \item The \texttt{event} keyword builds the observer pattern into the
      language (subscribe with \texttt{+=}, unsubscribe with \texttt{-=}) --
      Java has no direct counterpart.
  \end{itemize}
\end{frame}

\begin{frame}[fragile]{Delegates, Lambdas, and Events: Example}
  \framelabel{slide:alarm}
  \begin{lstlisting}
delegate void Notify(string msg);

class Alarm {
  public event Notify Triggered;

  public void Fire() {
    Triggered?.Invoke("Alarm!");
  }
}

Alarm a = new Alarm();
a.Triggered += msg => Console.WriteLine(msg);
a.Fire();
  \end{lstlisting}
  \hint{\texttt{event} hides the subscriber list, exposing only
    \texttt{+=}/\texttt{-=} to subscribe or unsubscribe.}
\end{frame}

\begin{frame}[fragile]{Action and Func}
  \begin{itemize}
    \item \texttt{Action<T1,...>} is a predefined generic delegate for
      a method returning \texttt{void}.
    \item \texttt{Func<T1,...,TResult>} is a predefined generic delegate
      for a method returning \texttt{TResult} (the last type parameter).
    \item Both come from the framework -- no need to declare your own
      \texttt{delegate} type for common cases.
  \end{itemize}
  \begin{lstlisting}
Action<string> print = msg => Console.WriteLine(msg);
print("Hello");

Func<int, int, int> add = (a, b) => a + b;
int sum = add(3, 4); // 7

Func<int, bool> isEven = n => n % 2 == 0;
List<int> numbers = new List<int> { 1, 2, 3, 4 };
List<int> evens = numbers.Where(isEven).ToList();
  \end{lstlisting}
  \hint{\texttt{Predicate<T>} is just a named alias for
    \texttt{Func<T, bool>}, used e.g.\ by \texttt{List<T>.Find}.}
\end{frame}

\begin{frame}[fragile]{Action and Func as Parameters}
  \texttt{Action}/\texttt{Func} are especially handy as
  \emph{parameter types} -- a method can accept a piece of behavior
  without declaring a custom delegate or interface for it.
  \begin{lstlisting}[basicstyle=\scriptsize\ttfamily]
static void Process(List<int> items, Action<int> onEach) {
  foreach (int item in items) onEach(item);
}

static int Reduce(List<int> items, int seed,
                   Func<int, int, int> combine) {
  int result = seed;
  foreach (int item in items) result = combine(result, item);
  return result;
}

List<int> numbers = new List<int> { 1, 2, 3, 4 };
Process(numbers, n => Console.WriteLine(n));
int sum = Reduce(numbers, 0, (a, b) => a + b); // 10
  \end{lstlisting}
  \hint{Unity APIs do this a lot, e.g.\ \texttt{button.onClick.AddListener(Action)}.}
\end{frame}

\begin{exerciseframe}{Events Are Not Fields}
  Take \texttt{Alarm} from \slideref{slide:alarm}. Subscribe two handlers,
  unsubscribe one of them, fire, and predict the output.
  \vspace{1ex}

  Then try these statements from \emph{outside} \texttt{Alarm} and read the error:
  \vspace{1ex}

  \begin{lstlisting}[basicstyle=\scriptsize\ttfamily]
Alarm a = new Alarm();
a.Triggered = null;
a.Triggered.Invoke("x");
  \end{lstlisting}
  \vspace{1ex}

  Finally replace \texttt{event Notify Triggered} by a plain field
  \texttt{Notify Triggered}. Which of the two statements compiles now?
\end{exerciseframe}

\begin{solutionframe}{Events Are Not Fields}
  \begin{itemize}
    \item Only the handler still subscribed runs, as \texttt{-=} removes
      exactly the one handler passed to it.
    \item Both statements referring to \texttt{Triggered} are rejected:\\
     \texttt{CS0070: The event 'Alarm.Triggered' can only appear on the left hand side of += or -=
      (except when used from within the type 'Alarm').}
    \item With a plain delegate field both compile -- any caller could then
      discard all subscribers or raise the alarm itself.
  \end{itemize}
  \vspace{3ex}

  \hint{That is what \texttt{event} buys you: outside the declaring type only
  \texttt{+=} and \texttt{-=} are permitted. Inside, \texttt{Triggered} is
  \texttt{null} while nobody is subscribed, which is why
  \texttt{Triggered?.Invoke(\ldots)} uses the null-conditional operator.}
\end{solutionframe}

\begin{exerciseframe}{One Method Instead of Two}

  Given the following two methods, which differ in their comparison only:
  \vspace{0.5ex}

  \begin{columns}[T]
    \column{0.5\textwidth}
      \begin{lstlisting}[basicstyle=\scriptsize\ttfamily]
static int SumOfEvens(List<int> items)
{
  int result = 0;
  foreach (int item in items)
  {
    if (item % 2 == 0)
    {
      result += item;
    }
  }
  return result;
}
      \end{lstlisting}
    \column{0.5\textwidth}
      \begin{lstlisting}[basicstyle=\scriptsize\ttfamily]
static int SumOfOdds(List<int> items)
{
  int result = 0;
  foreach (int item in items)
  {
    if (item % 2 != 0)
    {
      result += item;
    }
  }
  return result;
}
      \end{lstlisting}
  \end{columns}

  Replace both by a single method that receives the criterion as a parameter.

\end{exerciseframe}

\begin{solutionframe}{One Method Instead of Two}
  \begin{lstlisting}[basicstyle=\scriptsize\ttfamily]
static int Sum(List<int> items, Func<int, bool> keep)
{
  int result = 0;
  foreach (int item in items)
  {
    if (keep(item))
    {
      result += item;
    }
  }
  return result;
}

int evens = Sum(numbers, n => n % 2 == 0);   // 6
int odds = Sum(numbers, n => n % 2 != 0);    // 4
  \end{lstlisting}
  \vspace{2ex}

  \hint{Passing behavior instead of duplicating loops is the very shape behind
  LINQ's \texttt{Where}.}
\end{solutionframe}

\begin{exerciseframe}{Behavior as a Parameter}
  Implement \texttt{Retry} so that it calls \texttt{attempt} at most
  \texttt{times} times and yields \texttt{true} as soon as one call succeeds.
  \vspace{1ex}

  \begin{lstlisting}[basicstyle=\scriptsize\ttfamily]
static bool Retry(int times, Func<bool> attempt) { ... }
  \end{lstlisting}

  What do the following statements print?
  \vspace{1ex}

\begin{lstlisting}[basicstyle=\scriptsize\ttfamily]
int calls = 0;
Func<bool> succeedsOnThird = () => ++calls == 3;

Console.WriteLine(Retry(5, succeedsOnThird));
calls = 0;
Console.WriteLine(Retry(2, succeedsOnThird));
  \end{lstlisting}
\end{exerciseframe}

\begin{solutionframe}{Behavior as a Parameter}
  \begin{lstlisting}[basicstyle=\scriptsize\ttfamily]
static bool Retry(int times, Func<bool> attempt)
{
  for (int i = 0; i < times; i++)
  {
    if (attempt())
    {
      return true;
    }
  }
  return false;
}
  \end{lstlisting}
  The statements print \texttt{True} and \texttt{False}: three attempts are
  granted in the first case, only two in the second.
  \vspace{3ex}

  \hint{\texttt{Func<bool>} takes no arguments and returns a \texttt{bool}. Were
  there nothing to return, \texttt{Action} would be the counterpart.}
\end{solutionframe}

\begin{frame}[fragile]{Operator Overloading}
  \framelabel{slide:operators}
  \begin{itemize}
    \item Java: not allowed (the \texttt{+} operator on \texttt{String}
      is a special case built into the language).
    \item C\#: user-defined types can overload operators -- used
      extensively by Unity's math types.
  \end{itemize}
  \begin{lstlisting}
public static Vector3 operator +(Vector3 a, Vector3 b)
  => new Vector3(a.x + b.x, a.y + b.y, a.z + b.z);
  \end{lstlisting}
  \pause
  \vfill

  \hint{The \texttt{==} operator for \texttt{GameObject} in Unity is overloaded.
  When you evaluate\\
  myGameObject == null, Unity isn't just checking if the C\# reference
  is empty. It is checking if the underlying native C++ engine object has been
  destroyed.\\
  The C++ object is destroyed only at the end of the frame.}

\end{frame}

\begin{exerciseframe}{Operators on a Money Struct}
  Give \texttt{Money} an \texttt{operator +} adding two amounts and an
  \texttt{operator ==} letting equal amounts compare equal. Compile after each
  of the two steps and read what the compiler says about \texttt{==}.
  \vspace{1ex}

  \begin{lstlisting}[basicstyle=\scriptsize\ttfamily]
struct Money
{
  public int Cents;

  public Money(int cents)
  {
    Cents = cents;
  }
}

Money a = new Money(150);
Money b = new Money(99);
Console.WriteLine((a + b).Cents);        // 249
Console.WriteLine(a == new Money(150));  // True
  \end{lstlisting}
\end{exerciseframe}

\begin{solutionframe}{Operators on a Money Struct}
  \begin{lstlisting}[basicstyle=\scriptsize\ttfamily]
public static Money operator +(Money a, Money b)
  => new Money(a.Cents + b.Cents);

public static bool operator ==(Money a, Money b)
  => a.Cents == b.Cents;

public static bool operator !=(Money a, Money b)
  => !(a == b);
  \end{lstlisting}
  \begin{itemize}
    \item \texttt{==} alone does not compile: \texttt{CS0216} demands a matching
      \texttt{!=}.
    \item Two warnings remain: \texttt{CS0660} and \texttt{CS0661} ask you to
      override \texttt{Equals} and \texttt{GetHashCode} as well.
  \end{itemize}
  \vspace{3ex}

  \hint{Collections and LINQ compare with \texttt{Equals}, not with
  \texttt{==}, hence a type whose two notions of equality disagree behaves
  inconsistently. Unity's \texttt{GameObject} shows the other side of
  overloading \texttt{==}, see \slideref{slide:operators}.}
\end{solutionframe}

\begin{frame}{Exceptions}
  \begin{itemize}
    \item Java distinguishes \emph{checked} exceptions (must be declared
      with \texttt{throws} or caught) from \emph{unchecked} ones
      (\texttt{RuntimeException} and subclasses).
    \item C\# has only unchecked exceptions -- no \texttt{throws}
      declaration exists, and the compiler never forces you to catch
      anything.
  \end{itemize}
\end{frame}

\begin{exerciseframe}{A Signature Without throws}
  Port this method to C\#. What becomes of the \texttt{throws} clause, and what
  happens if the caller leaves out the \texttt{try}?
  \vspace{1ex}

  \begin{lstlisting}[language=Java,basicstyle=\scriptsize\ttfamily]
static String readFirstLine(String path)
    throws IOException
{
  ...
}

try {
  System.out.println(readFirstLine("in.txt"));
} catch (IOException e) {
  System.out.println("cannot read");
}
  \end{lstlisting}
  \vspace{1ex}

  How does a caller of the C\# version learn that reading may fail?
\end{exerciseframe}

\begin{solutionframe}{A Signature Without throws}
  \begin{itemize}
    \item The clause simply disappears, as no \texttt{throws} exists in C\#:
      \texttt{static string ReadFirstLine(string path)}.
    \item Leaving out the \texttt{try} still compiles. Nothing reminds the
      caller, the exception propagates and terminates the program. Unity logs it
      and aborts the current callback.
    \item Only documentation remains: possible exceptions are recorded in XML
      documentation comments, which the IDE shows at the call site.
  \end{itemize}
  \vspace{3ex}

  \hint{\texttt{try}, \texttt{catch} and \texttt{finally} work as in Java. What
  is missing is the compiler forcing you to handle anything, which moves the
  burden to documentation and reviews.}
\end{solutionframe}

\begin{frame}[fragile]{Collections: Streams vs.\ LINQ}
  \framelabel{slide:linq}
  \begin{columns}[T]
    \column{0.5\textwidth}
      \textbf{Java Streams}
      \begin{lstlisting}[language=Java]
names.stream()
  .filter(n -> n.length() > 3)
  .map(String::toUpperCase)
  .collect(Collectors.toList());
      \end{lstlisting}
    \column{0.5\textwidth}
      \textbf{C\# LINQ}
      \begin{lstlisting}
names
  .Where(n => n.Length > 3)
  .Select(n => n.ToUpper())
  .ToList();
      \end{lstlisting}
  \end{columns}
\end{frame}

\begin{exerciseframe}{When Does a LINQ Query Run?}
  The pipeline of \slideref{slide:linq} looked like a single computation.
  Predict the two numbers returned by \texttt{Count()} here.
  \vspace{1ex}

  \begin{lstlisting}[basicstyle=\scriptsize\ttfamily]
List<string> names = new List<string> { "Ann", "Bernard" };
IEnumerable<string> query = names.Where(n => n.Length > 3);
names.Add("Christina");

Console.WriteLine(query.Count());
Console.WriteLine(query.Count());
  \end{lstlisting}
  \vspace{1ex}

  What would the second call do in Java, applied to a stream?
\end{exerciseframe}

\begin{solutionframe}{When Does a LINQ Query Run?}
  \begin{itemize}
    \item Both print \texttt{2}. The query runs only when enumerated, hence it
      sees \texttt{"Christina"}, which was added after the query was built.
    \item It may be enumerated as often as you like. A Java stream may be
      consumed once only, a second terminal operation throws an
      \texttt{IllegalStateException}.
    \item Both \texttt{ToList()} and \texttt{ToArray()} run the query once and
      freeze the outcome.
  \end{itemize}
  \vspace{5ex}

  \hint{Deferred execution is what lets a query read like a description rather
  than a computation. Beware of sources changing in between, and of enumerating
  an expensive query more than once.}
\end{solutionframe}

\begin{frame}[fragile]{Small Syntax Differences}
  \begin{itemize}
    \item String interpolation: \texttt{\$"Hello, \{name\}!"} in C\#
      vs.\ \texttt{"Hello, " + name + "!"} or \texttt{String.format(...)}
      in Java.
    \item Both languages have \texttt{var} for local type inference; C\#
      additionally allows target-typed \texttt{new()}, e.g.\
      \texttt{List<int> xs = new();}
    \item C\# supports \emph{extension methods} -- adding a method to an
      existing type without subclassing or modifying it. Java has no
      equivalent (default interface methods come closest but require an
      interface).
    \item Generics: C\# generics are reified (the type is known at run
      time); Java generics use type erasure.
    \item \texttt{nameof(x)} yields the name of \texttt{x} as a string, hence
      parameter names passed to exceptions survive a rename. Java has no
      equivalent.
  \end{itemize}
  \vfill

  \hint{Do not use \texttt{var}. You would suppress important type information.}

\end{frame}

\begin{frame}[fragile]{Extension Methods}
  \framelabel{slide:extensions}
  \begin{itemize}
    \item An extension method is a \texttt{static} method in a non-generic
      \texttt{static} class whose first parameter carries the \texttt{this}
      modifier, naming the type to be extended.
    \item It is called as if it were an instance method of that type; the
      compiler turns such a call into the plain static call.
  \end{itemize}
  \vspace{1ex}

  \begin{columns}[T]
    \column{0.55\textwidth}
      \textbf{Declaration}
      \begin{lstlisting}[basicstyle=\scriptsize\ttfamily]
public static class IntExtensions
{
  public static bool IsEven(this int n)
  {
    return n % 2 == 0;
  }
}
      \end{lstlisting}
    \column{0.45\textwidth}
      \textbf{Two calls, one meaning}
      \begin{lstlisting}[basicstyle=\scriptsize\ttfamily]
7.IsEven();

IntExtensions.IsEven(7);
      \end{lstlisting}
  \end{columns}
  \vspace{1ex}

  \begin{itemize}
    \item An extension method reaches only what is accessible from outside, and
      it loses against an instance method of the same signature.
    \item LINQ rests upon them: \texttt{Where} and \texttt{Select} are extension
      methods of \texttt{IEnumerable<T>}, which is why they chain onto any
      collection (\slideref{slide:linq}).
  \end{itemize}
\end{frame}

\begin{exerciseframe}{An Extension Method}
  Write \texttt{IsBlank} as an extension method on \texttt{string}, yielding
  \texttt{true} for the empty string and for a string of blanks. Where does the
  compiler insist that the method be declared?
  \vspace{1ex}

  \begin{lstlisting}[basicstyle=\scriptsize\ttfamily]
Console.WriteLine("   ".IsBlank());   // True
Console.WriteLine("x".IsBlank());     // False

string missing = null;
Console.WriteLine(missing.IsBlank()); // ?
  \end{lstlisting}
  \vspace{1ex}

  Predict the last line, then compare it with what a normal instance method
  would do.
\end{exerciseframe}

\begin{solutionframe}{An Extension Method}
  \begin{lstlisting}[basicstyle=\scriptsize\ttfamily]
public static class StringExtensions
{
  public static bool IsBlank(this string s)
    => s == null || s.Trim().Length == 0;
}
  \end{lstlisting}
  \begin{itemize}
    \item The enclosing class has to be static and non-generic, otherwise
      \texttt{CS1106} results.
    \item The last line prints \texttt{True} without throwing: an extension
      method is a static call, hence \texttt{s} may well be \texttt{null}. A
      genuine instance method would throw a \texttt{NullReferenceException}.
  \end{itemize}
  \vspace{1ex}

  \hint{That is the essential difference from a real method: the receiver is
  just the first argument, and there is no dispatch on it. Java has no
  equivalent.}
\end{solutionframe}

\begin{frame}{Extension Methods: Resolution}
  Given both \texttt{Describe(this Base)} and \texttt{Describe(this Derived)}
  where \texttt{Derived} inherits from \texttt{Base},
  the \emph{static} type of the receiver decides. There is no virtual dispatch.
  \vspace{1ex}

  \begin{center}
    \footnotesize
    \begin{tabular}{ll}
      \hline
      \textbf{Expression} & \textbf{Method selected}\\
      \hline
      \texttt{Base b = new Derived(); b.Describe()} & extension on \texttt{Base}\\
      \texttt{new Derived().Describe()}             & extension on \texttt{Derived}\\
      \texttt{((Derived)b).Describe()}              & extension on \texttt{Derived}\\
      \hline
    \end{tabular}
  \end{center}
  \vspace{1ex}

  \begin{itemize}
    \item An instance method of matching signature always wins, an inherited one
      as well. Adding a method to a base class can hence silently take calls
      away from an extension method.
    \item An extension method in a nearer namespace shadows one further out,
      which is as close as C\# gets to overriding an extension method.
    \item Two applicable extension methods at the same distance are ambiguous
      (\texttt{CS0121}).
  \end{itemize}
\end{frame}

\begin{exerciseframe}{Which Method Runs?}
  Both extension methods are in scope. Predict the three outputs before you run
  the program.
  \vspace{1ex}

  \begin{lstlisting}[basicstyle=\scriptsize\ttfamily]
class Base { }
class Derived : Base { }

static class BaseExtensions
{
  public static string Describe(this Base b) => "Base";
  public static string Describe(this Derived d) => "Derived";
}

Base b = new Derived();
Console.WriteLine(b.Describe());
Console.WriteLine(new Derived().Describe());
Console.WriteLine(((Derived)b).Describe());
  \end{lstlisting}
\end{exerciseframe}

\begin{solutionframe}{Which Method Runs?}
  \begin{lstlisting}[language={},basicstyle=\scriptsize\ttfamily]
Base
Derived
Derived
  \end{lstlisting}
  \begin{itemize}
    \item The first line is the surprising one: the object is a
      \texttt{Derived}, yet the extension of \texttt{Base} runs, because
      \texttt{b} is \emph{declared} as \texttt{Base}.
    \item The cast in the third line changes the outcome, which would never
      happen with a virtual instance method.
  \end{itemize}
  \vspace{4ex}

  \hint{Resolution takes place at compile time on the static type. Extension
  methods are therefore a poor fit wherever polymorphic behavior is wanted.}
\end{solutionframe}

\begin{frame}[fragile]{Virtual and Override}
  \framelabel{slide:virtual}
  \begin{itemize}
    \item Java: every method is virtual, and \texttt{@Override} is an optional
      annotation letting the compiler check your intention.
    \item C\#: a method is \emph{not} virtual by default. The base class has to
      offer \texttt{virtual} (or \texttt{abstract}), and the derived class has to
      state \texttt{override}.
    \item Overriding something not virtual is an error (\texttt{CS0506}).
      Repeating a signature without \texttt{override} hides the inherited method
      and warns (\texttt{CS0108}), while \texttt{new} declares the hiding as
      intended.
  \end{itemize}
  \vspace{1ex}

  \begin{lstlisting}[basicstyle=\scriptsize\ttfamily]
class Base
{
  public virtual string Speak() => "Base";
  public string Walk() => "Base";
}

class Sub : Base
{
  public override string Speak() => "Sub";  // overrides
  public new string Walk() => "Sub";        // hides
}
  \end{lstlisting}
\end{frame}

\begin{exerciseframe}{Override or Hide?}
  Predict the three outputs for the classes of \slideref{slide:virtual}, then
  explain the difference.
  \vspace{1ex}

  \begin{lstlisting}[basicstyle=\scriptsize\ttfamily]
Base asBase = new Sub();
Console.WriteLine(asBase.Speak());
Console.WriteLine(asBase.Walk());
Console.WriteLine(((Sub)asBase).Walk());
  \end{lstlisting}
  \vspace{1ex}

  And what happens if you drop \texttt{new} from \texttt{Sub.Walk}, or write
  \texttt{override} there instead?
\end{exerciseframe}

\begin{solutionframe}{Override or Hide?}
  \begin{lstlisting}[language={},basicstyle=\scriptsize\ttfamily]
Sub
Base
Sub
  \end{lstlisting}
  \begin{itemize}
    \item \texttt{Speak} is virtual and overridden, hence the run-time type
      decides.
    \item \texttt{Walk} is hidden rather than overridden, hence the static type
      decides, and only the cast reaches \texttt{Sub.Walk}.
    \item Without \texttt{new} the code still compiles, but warns
      (\texttt{CS0108}). With \texttt{override} it no longer compiles
      (\texttt{CS0506}), as \texttt{Base.Walk} is not virtual.
  \end{itemize}
  \vspace{4ex}

  \hint{In Java the second line would print \texttt{Sub}, since every instance
  method is virtual there and hiding one is not possible at all.}
\end{solutionframe}

\subsection{Generics}

\begin{frame}[fragile]{Generics Are Reified, Not Erased}
  \framelabel{slide:generics}
  \begin{lstlisting}[basicstyle=\scriptsize\ttfamily]
List<int> numbers = new();   // really ints: no boxing, no Integer

Type t = typeof(T);          // the actual type argument, at run time
bool b = x is T;             // works
T fresh = new T();           // legal -- with a new() constraint
T[] array = new T[10];       // legal
  \end{lstlisting}
  \vspace{1ex}

  \begin{itemize}
    \item Java erases its type arguments: \texttt{List<String>} and
      \texttt{List<Integer>} are one class at run time, \texttt{new T()} is
      impossible, and a list of \texttt{int} boxes every element.
    \item C\# keeps them. The runtime compiles a separate body for every value
      type and shares one among all reference types, so \texttt{List<int>}
      stores plain \texttt{int}s.
    \item Because the type survives, \texttt{typeof(T)}, \texttt{is~T},
      \texttt{default(T)} and \texttt{new T()} all mean something --- which is
      why the reflection on \slideref{slide:typeobject} can work with generic
      code at all.
  \end{itemize}
\end{frame}

\begin{frame}[fragile]{Constraints: What You Promise About \texttt{T}}
  \begin{lstlisting}[basicstyle=\scriptsize\ttfamily]
where T : GraphElement              // a base class: use its members
where T : Component                 // Unity's Component
where T : IPersistentConfigItem, new()  // an interface, and constructible
where E : struct, IConvertible      // a value type -- the old way to say enum
where T : class                     // any reference type
  \end{lstlisting}
  \vspace{1ex}

  \begin{itemize}
    \item Without a constraint a \texttt{T} is no more than an
      \texttt{object}: you may pass it around and call
      \texttt{ToString()} on it, and that is all.
    \item Every constraint buys one thing. A base class or interface gives you
      its members, \texttt{new()} lets you construct a \texttt{T},
      \texttt{struct} and \texttt{class} settle boxing and nullability.
    \item \texttt{where E : struct, IConvertible} is the idiom from before
      C\# 7.3, still used in \texttt{ConfigIO.RestoreEnum}. Today one writes
      \texttt{where E : Enum} and means it exactly.
  \end{itemize}
\end{frame}

\begin{frame}[fragile]{One Algorithm, Two Element Types}
  \begin{lstlisting}[basicstyle=\scriptsize\ttfamily]
// Assets/SEE/DataModel/DG/MergeDiffGraphExtension.cs
private static void MergeGraphElements<T>(ISet<T> added, ISet<T> removed,
        ISet<T> changed, Action<T> addToGraph, Graph baseline)
    where T : GraphElement
{
    added.ForEach(e => e.SetToggle(ChangeMarkers.IsNew));   // needs the
    ...                                                     // constraint
}
...
MergeGraphElements(addedNodes, removedNodes, changedNodes, ...);  // nodes
MergeGraphElements(addedEdges, removedEdges, changedEdges, ...);  // edges
  \end{lstlisting}
  \vspace{1ex}

  \begin{itemize}
    \item Drop the \texttt{where} and \texttt{SetToggle} stops compiling. The
      constraint is what turns a method that shuffles sets into a graph
      algorithm.
    \item One body serves nodes and edges, and each call keeps its own type:
      the \texttt{Action<T>} that adds to the graph is
      \texttt{AddNode} in one case and \texttt{AddEdge} in the other, checked
      at compile time.
    \item The alternative --- parameters of type \texttt{GraphElement} and a
      cast at every use --- gives up exactly that check.
  \end{itemize}
\end{frame}

\begin{exerciseframe}{Make It Generic}
  \begin{enumerate}
    \item Two methods differ only in their element type: one returns the
      \texttt{Node} with the largest value of a metric, the other the
      \texttt{Edge}. Merge them into one method. Which constraint do you need,
      and what breaks without it?
    \item \texttt{ConfigIO.RestoreList<T>} is declared \texttt{where T :
      IPersistentConfigItem, new()} and its body says \texttt{T t = new();}.
      It could have called \texttt{Activator.CreateInstance} instead --- why is
      the constraint better?
    \item \texttt{where E : struct, IConvertible} does not actually restrict
      \texttt{E} to enumerations. Name a type that satisfies it and is not an
      enum, and say what you would write today.
  \end{enumerate}
\end{exerciseframe}

\begin{solutionframe}{Make It Generic}
  \begin{enumerate}
    \item \texttt{where T : GraphElement}, because the body has to read a
      metric --- \texttt{TryGetNumeric} is declared on \texttt{Attributable},
      inherited by \texttt{GraphElement}. Without the constraint \texttt{T} is
      an \texttt{object} and no member of it is in scope.
    \item Both construct a \texttt{T}, but \texttt{new()} is checked when the
      code is compiled, and by the \emph{caller}: a type without a public
      parameterless constructor cannot even be passed.
      \texttt{Activator.CreateInstance} moves the same requirement to run time,
      where it becomes an exception --- the trade-off of
      \slideref{slide:typeobject} in miniature.
    \item Any value type implementing \texttt{IConvertible} qualifies:
      \texttt{int}, \texttt{double}, \texttt{bool}, \texttt{DateTime}. Since
      C\# 7.3 the honest spelling is \texttt{where E : Enum}.
  \end{enumerate}
\end{solutionframe}

%%% Local Variables:
%%% mode: latex
%%% TeX-master: "main-slides"
%%% mode: reftex
%%% mode: flyspell
%%% End:
